% !TEX program = xelatex
% 编译方式：在本目录执行两遍 xelatex（详见 docs/_manual_common/README.md）。
% 源码依据：src/analysis/ 与 include/hacdcpf/analysis/
%   （hosting_capacity / multidimensional_weak_link / counterfactual_planning），公式与参数以代码为准。
\documentclass[UTF8,fontset=fandol,a4paper,11pt]{ctexart}
\usepackage{../../_manual_common/hysim_manual}

\title{\textbf{交直流混合高弹性能源电力系统仿真平台}\\[0.5em]
       {\Large 承载力与薄弱环节分析技术手册}\\[0.3em]
       {\large 数学模型、约束、算法流程、数值稳定性与工业级评估}}
\author{HySim-XJTU-HRPES（西安交通大学 HRPES 团队）}
\date{\today}

\begin{document}
\maketitle
\begin{abstract}
\noindent 本手册依据 HySim-XJTU-HRPES 平台分析模块
（\srcpath{src/analysis/}、\srcpath{include/hacdcpf/analysis/}）整理，覆盖三类面向
规划与评估的后处理方法：依据《DL/T 2041-2025》的设备级分布式资源承载力评估、时空
多维度薄弱环节筛查、反事实规划措施评估。全部结论以源码为准：承载力工程校验默认关闭
（仅做规划筛查），薄弱环节为非补偿式多准则筛查与决策路由（不宣称归因量等于反事实投资
收益、缺失维度不补零），反事实措施施加于系统副本、经济/可靠性/弹性指标使用选项给定
的有界预算与时窗。命名空间为 \texttt{hacdcpf::analysis}。文档风格、参数表体例与符号
约定与《潮流计算技术手册》《最优潮流技术手册》完全一致，全部近似、回退与模型覆盖
边界均如实标注。
\end{abstract}

\tableofcontents
\newpage

\section{阅读指南与约定}
\subsection{数据来源}
本手册全部内容取自以下源代码与既有契约文档，公式与参数均以代码为准：
\begin{itemize}
  \item \srcpath{include/hacdcpf/analysis/hosting\_capacity.hpp}、
        \srcpath{src/analysis/hosting\_capacity.cpp} —— 承载力评估接口与实现；
  \item \srcpath{include/hacdcpf/analysis/multidimensional\_weak\_link.hpp}、
        \srcpath{src/analysis/multidimensional\_weak\_link.cpp} —— 多维薄弱环节筛查；
  \item \srcpath{include/hacdcpf/analysis/counterfactual\_planning.hpp}、
        \srcpath{src/analysis/counterfactual\_planning.cpp} —— 反事实规划措施评估；
  \item \srcpath{docs/theory/capacity\_analysis\_implementation.md}、
        \srcpath{docs/theory/multidimensional\_weak\_link\_identification.md} —— 既有契约文档。
\end{itemize}
上述头文件与 \srcpath{include/hacdcpf/analysis/} 同目录，其余同名头文件
（\srcpath{short\_circuit.hpp}、\srcpath{scenario\_generation.hpp} 等）实现于兄弟模块，
不属于本手册范围。文中“\texttt{文件:行号}”引用基于当前检出，代码演进后行号可能漂移，
函数名与文件路径保持稳定。

\subsection{单位制与索引空间}
承载力 $S_d$ 单位 MVA/MW，功率因数 $\cos\theta$ 无量纲；薄弱环节压力量归一化后无量纲，
时间维度以选项声明的评估周期为索引空间；反事实四维收益向量的物理量分别为经济运行成本
（\$/yr）、碳排放（tCO\textsubscript{2}/yr）、可靠性 EENS（MWh/yr）与弹性加权失负荷 ENS
（MWh）。承载力评估以变压器下游供电区（BFS 展开）为索引，结果按 \fld{transformers} 与
\fld{areas} 分列；区域聚合以 \fld{bus.area}（县域）为索引空间。

\subsection{符号约定}
\begin{itemize}
  \item $S$：变压器额定容量（MVA），并联时乘 \fld{n\_parallel}；$P$：供电区聚合负荷；
        $P_G$：区内常规（非分布式）发电；$\beta$：最大反向负载率；$\tau_{\max}$：最大出力
        系数；$P_{\text{ESS}},\Delta P_{\text{ESS}}$：现有/预期储能充电功率。
  \item $x_{ik}$：实体 $i$ 在维度 $k$ 的原始压力证据；$\tau_k$：维度 $k$ 的工程目标阈值；
        $p_{ik}$：无量纲越限比；$\sigma_i,c_i,d_i$：瓶颈严重度、共识度、维度分歧。
  \item $\mathbf{m}=(C_{\text{op}},E_{\mathrm{CO_2}},\text{EENS},\text{ENS})^\top$：
        反事实四维度量向量；$\Delta\mathbf{m}$：相对基线收益。
\end{itemize}

\subsection{诚实口径}
三类分析均为规划/评估后处理，不改变调度：承载力 \fld{enable\_verification} 默认为
\texttt{false}，即默认只给规划筛查值、不回跑潮流/短路/谐波引擎校核；薄弱环节为非补偿式
筛查，缺失维度不被当作零、也不宣称归因排放/失负荷等于反事实投资收益；反事实措施施加于
系统副本，经济/可靠性/弹性指标使用 \fld{reliability\_physical\_budget}、
\fld{resilience\_horizon\_hours} 等有界预算与时窗。相关限制在第~\ref{sec:limits}~节集中列出。

\section{模块定位与工程应用场景}\label{sec:scope}
\subsection{功能定位}
分析模块位于求解引擎之上、GUI 之下，是\textbf{面向规划与评估的后处理层}：它调用潮流、
短路、谐波、可靠性、弹性、碳流等引擎产生的证据，输出可供决策的分级、排序与措施收益，
但\textbf{不}参与实时调度，也\textbf{不}回写系统状态。三条子线各自独立：
\begin{enumerate}
  \item \textbf{承载力评估}：回答“某台配电变压器下游还能接入多少分布式资源”，按
        《DL/T 2041-2025》给出设备级绿/黄/红分级与可选工程校验；
  \item \textbf{多维薄弱环节筛查}：在经济、碳、可靠性、弹性四个维度上做非补偿式脆弱性
        筛查，定位需优先加固/改造的空间实体与关键时段；
  \item \textbf{反事实规划}：对候选工程措施（增容、储能、联络开关、自动化、DER）在系统
        副本上逐条评估四维收益并度量两两协同，为投资排序提供依据。
\end{enumerate}

\subsection{工程应用场景}
典型场景包括：配电网分布式光伏/储能接入的规划审查（承载力）；县域—馈线两级的迎峰度夏
薄弱点排查（薄弱环节 \texttt{operation} 模式）；五年规划投资组合的措施比选与协同分析
（反事实 \texttt{planning} 模式）。三者可串联使用：薄弱环节筛查定位问题实体 $\to$ 反事实
评估比选措施 $\to$ 承载力评估复核接入余量。

\section{输入、输出与边界条件}
\subsection{输入}
\begin{itemize}
  \item \textbf{公共输入}：已装配的 \texttt{HybridPowerSystem}（工程语义层）与对应选项结构体；
        承载力另需变压器/储能上的 \fld{cap\_*} 前缀属性（随算例往返保存）。
  \item \textbf{承载力}：全局默认参数（$\cos\theta,\tau_{\max},\beta$、N-1 上限、工程校验开关、
        电压偏差限、THD 限）；元件级 \fld{cap\_*} 属性优先于全局默认。
  \item \textbf{薄弱环节}：来自经济/碳/可靠性/弹性四引擎的原始压力证据 $x_{ik}$、阈值 $\tau_k$、
        决策模式（\texttt{planning}/\texttt{operation}）、时段集合。
  \item \textbf{反事实}：措施族启用开关、预算/时窗上界、能量折价。
\end{itemize}

\subsection{输出}
承载力输出 \fld{transformers[]}（每台 $S_d,C_{d1},C_{d2}$ 区间与分级）、\fld{areas[]}
（县域聚合与分级）、\fld{warnings[]}、\fld{verification\{\}}；薄弱环节输出实体排序、
关键时段、Pareto 前沿计数与冲突/兼容计数；反事实输出基线度量、逐措施收益与两两协同项。
全部结果可 JSON 序列化并经 GUI 渲染。

\subsection{边界条件}
承载力要求供电区可由“同压支路邻接 + 变压器边界”BFS 明确划分，孤立/失电母线不计入下游区；
薄弱环节要求实体具有相同的可优化组件组与相同覆盖集方可 Pareto 比较，缺失维度落在 $K_i$ 之外；
反事实要求措施可在系统副本上物理施加且引擎可解，不可解措施记为无效并跳过。

\section{变量、参数与符号定义}
\subsection{关键数据结构}
\compmeta{HostingCapacityResult}{include/hacdcpf/analysis/hosting\_capacity.hpp}
主要字段：\fld{standard}（默认 \texttt{"DL/T 2041-2025"}）、\fld{transformers}
（\texttt{TransformerHostingResult} 数组）、\fld{areas}（\texttt{AreaHostingResult} 数组）、
\fld{warnings}、\fld{verification}（\texttt{HostingVerification}）。

\compmeta{MultidimensionalWeakLinkResult}{include/hacdcpf/analysis/multidimensional\_weak\_link.hpp}
主要字段：\fld{schema}、\fld{dimension\_coverage}（四维覆盖）、\fld{entities}、
\fld{critical\_periods}、\fld{pareto\_count}、\fld{compatible\_count}、\fld{conflict\_count}。

\compmeta{CounterfactualPlanningResult}{include/hacdcpf/analysis/counterfactual\_planning.hpp}
主要字段：\fld{schema}、\fld{baseline}（\texttt{CounterfactualMetricSet}）、
\fld{measures}、\fld{interactions}、\fld{methodology}；度量向量为经济运行成本、碳、可靠性
EENS、弹性加权 ENS 四维。

\subsection{参数字典}
\begin{paramtable}
\fld{default\_power\_factor} & $\cos\theta$ & — & $0\sim1$ & $0.95$ & 承载力默认功率因数\\
\fld{default\_dr\_max\_output\_coeff} & $\tau_{\max}$ & — & $\ge10^{-9}$ & $1.0$ & 分布式电源最大出力系数默认\\
\fld{single\_transformer\_beta} & $\beta$ & — & $\ge0$ & $0.80$ & 单变压器可用容量系数\\
\fld{n1\_loading\_limit} & — & — & $\ge0$ & $1.0$ & N-1 负载率上限\\
\fld{enable\_verification} & — & 布尔 & — & \texttt{false} & 是否启用 §12 引擎校核\\
\fld{thd\_limit\_pct} & — & \% & — & $5.0$ & 谐波总畸变（THD）限值\\
\fld{top\_k} & $k$ & — & 整数 & $15$ & 输出前 $k$ 个薄弱环节\\
\fld{high\_pressure\_threshold} & $h$ & — & $0\sim1$ & $0.75$ & 高压力维度判定阈值\\
\fld{compatibility\_threshold} & $c$ & — & $0\sim1$ & $0.50$ & 措施兼容判定阈值\\
\fld{conflict\_gap\_threshold} & $g$ & — & $0\sim1$ & $0.60$ & 冲突间隙判定阈值\\
\fld{max\_measures} & — & — & 整数 & $5$ & 反事实候选措施上限\\
\fld{reliability\_physical\_budget} & — & — & 整数 & $40$ & 可靠性评估物理预算\\
\fld{resilience\_horizon\_hours} & — & h & — & $12$ & 弹性评估时窗\\
\fld{energy\_price\_per\_mwh} & — & \$/MWh & — & $100$ & 失负荷/电量折价\\
\end{paramtable}
反事实措施开关 \fld{enable\_line\_capacity}、\fld{enable\_storage}、\fld{enable\_tie\_switch}、
\fld{enable\_automation}、\fld{enable\_der} 控制候选措施族的启用；薄弱环节 \fld{mode} 取
\texttt{Planning}/\texttt{Operation}，\fld{target\_pressure} 为四维目标压力向量。

\section{源码等价分析模型}

本章仅转写 \srcpath{src/analysis/hosting\_capacity.cpp}、
\srcpath{multidimensional\_weak\_link.cpp} 与
\srcpath{counterfactual\_planning.cpp} 当前实际执行的计算。规范名称只说明功能背景，
不用于补充源码中不存在的公式。

\subsection{承载力评估}

\srcpath{hosting\_capacity\_options\_from\_json} 对参数执行以下截断：功率因数与
$k_r$ 截断到 $[0,1]$，\fld{default\_dr\_max\_output\_coeff} 下限为 $10^{-9}$，
其余 $\beta$、N-1 负载限值、电压偏差和 THD 限值下限为零
（\srcpath{src/analysis/hosting\_capacity.cpp:hosting\_capacity\_options\_from\_json}）。

代码只沿两端额定电压差小于 $10^{-3}$ kV 的在运 ACBranch 建立无向邻接；其余支路均作为
变压器边界。由低压侧开始的 BFS 只访问在运母线，所得集合用于累加
\begin{align}
 P&=\sum_b\left[\max(0,P_{d,b})+sum_{l\in b}\max(0,P_l\,\fld{scaling}_l)\right],\\
 P_G&=\sum_{g\in b,\,on}P_g,\\
 P_{DR}&=\sum_{pv,rg\in b,\,on}P+
          \sum_{sg\in b,\,on}\max(0,P_{sg}\,\fld{scaling}_{sg}),\\
 P_{ESS}&=\sum_{s\in b,\,on}\begin{cases}
 \max(0,\fld{cap\_static\_charging\_mw}_s),&\fld{cap\_charging\_strategy}=\texttt{static},\\
 \max(0,-P_s),&\text{其他字符串}.
 \end{cases}
\end{align}
依据为 \srcpath{src/analysis/hosting\_capacity.cpp:assess\_hosting\_capacity}。

每台设备取 $n=\max(1,\fld{n\_parallel})$。设备功率因数和最大出力系数仅在设备字段大于零时
使用，否则采用选项默认值；若所得 $\tau\le0$，再强制为 1。若设备 $\beta>0$ 则直接使用，
否则
\begin{equation}
 \beta=\begin{cases}
 \fld{single\_transformer\_beta},&n\le1,\\
 \dfrac{n-1}{n}\fld{n1\_loading\_limit},&n>1.
 \end{cases}
\end{equation}
令 $S_{tot}=nS_N$、$P_{rev}=\beta S_{tot}pf$、
$B=P-P_G+P_{rev}+P_{ESS}$，代码分别计算
\begin{align}
 S_{d,min}&=\max\left(0,\frac{B+\Delta P_{ESS,min}}{\tau}\right),\\
 S_{d,max}&=\max\left(0,\frac{B+\Delta P_{ESS,max}}{\tau}\right),
\end{align}
若前者大于后者则交换。这里不额外限制 $S_N$、$pf$ 或 $\beta$ 的符号。
对应源码为 \srcpath{src/analysis/hosting\_capacity.cpp:pair\_key（lambda）}。

设备余量按赋值顺序为
\begin{align}
 C_{grid,min/max}&=S_{d,min/max}-P_{DR},\\
 C_{reg,min/max}&=C_{grid,min/max}-P_{registered}.
\end{align}
分级函数严格为：$C_{reg,min}>0$ 为 green；否则 $C_{grid,min}>0$ 为 yellow；否则 red。
区域先逐设备求和 $S_d$、已接入和已备案量，再重复同一赋值和分级。
最终设备等级只在所属区域为 red 时强制 red，区域 yellow 不覆盖设备自身等级。
依据为 \srcpath{src/analysis/hosting\_capacity.cpp:pair\_key（lambda）}。

启用工程校核时，潮流固定使用 \fld{max\_iter=100}、\fld{tol=$10^{-8}$}。
电压偏差为 $100(V_i-1)$，越过 $+\Delta U_H$ 或 $-\Delta U_L$ 才违规；
支路只使用 from 端
\begin{equation}
 loading=100\frac{\sqrt{P_f^2+Q_f^2}}{\fld{rate\_a\_mva}},
\end{equation}
且只在大于 100 时违规。短路批处理异常被吞掉并保留初始通过状态；谐波异常只把
\fld{harmonic\_evaluated} 设为假。推荐由违规字符串决定：存在 critical 为
\texttt{suspend\_connection}，否则有任意违规为 \texttt{allow\_with\_mitigation}，否则
\texttt{allow\_connection}。依据为 \srcpath{src/analysis/hosting\_capacity.cpp:pair\_key（lambda）}。

\subsection{多维薄弱环节}

对维度 $d\in\{0,1,2,3\}$，只有 pressure 存在且有限时才标记可用。令
\begin{equation}
 x_d=\max(0,\fld{pressure}_d),\qquad
 q_d=\frac{x_d}{\begin{cases}t_d,&t_d\text{ 有限且 }>0,\\1,&\text{否则}.
 \end{cases}}
\end{equation}
缺失维度的存储分数为零，但不参与证据数、最小值和最大值。
严重度是可用 $q_d$ 的最大值；共识度在有证据时为最小值，否则为零。
主导维度只在严格大于当前最大值时更新，因此所有可用分数均为零时可保持空字符串。
依据为 \srcpath{src/analysis/multidimensional\_weak\_link.cpp:normalized\_pressure}。

分歧度只在证据维数大于 1 时取 $\max q_d-\min q_d$，否则为零。关系按以下顺序判定，
先命中的分支结束：
\begin{enumerate}
 \item 证据数小于 $\max(1,\fld{minimum\_dimensions})$：\texttt{insufficient\_evidence}；
 \item 证据数至少 3 且共识度不小于 \fld{compatibility\_threshold}：\texttt{compatible}；
 \item 证据数至少 2、严重度不小于 \fld{high\_pressure\_threshold}、分歧度不小于
       \fld{conflict\_gap\_threshold}：\texttt{conflict}；
 \item 严重度不小于高压阈值：\texttt{single\_dimension}；否则 \texttt{watch}。
\end{enumerate}
依据为 \srcpath{src/analysis/multidimensional\_weak\_link.cpp:run\_multidimensional\_weak\_link\_assessment}。

实体 $a$ 支配实体 $b$ 仅当 comparison group 与四位覆盖掩码相同，并且对每个可用维度
\begin{equation}
 q_{a,d}+10^{-12}\ge q_{b,d},
\end{equation}
且至少一维满足 $q_{a,d}>q_{b,d}+10^{-12}$。证据不足的实体不进入 Pareto 判定。
排序先按 Pareto 真值，再在差值超过 $10^{-12}$ 时按严重度、共识度降序，最后按 key 升序；
之后才截断 top-k。依据为 \srcpath{src/analysis/multidimensional\_weak\_link.cpp:coverage\_mask} 和
\srcpath{src/analysis/multidimensional\_weak\_link.cpp:run\_multidimensional\_weak\_link\_assessment}。

\subsection{反事实规划指标和收益}

经济或碳指标启用时先以 \fld{max\_iter=100}、\fld{tol=$10^{-7}$} 求潮流。
经济指标优先取“收敛、目标有限且绝对值大于 $10^{-9}$”的 ACOPF 目标乘
\fld{annual\_hours}。否则在潮流可用时计算
\begin{equation}
 L=\sum_{branch}\max(0,P_f+P_t)
   +\sum_{vsc}\max(0,P_{loss})+\sum_{dcdc}\max(0,P_{loss}),
\end{equation}
并赋值 $L\cdot\fld{energy\_price\_per\_mwh}\cdot\fld{annual\_hours}$。
碳指标选择 matrix solved 时的 matrix summary，否则选择 tracing summary，并把
\fld{total\_generation\_emissions\_tco2} 乘 annual hours。
可靠性固定配置确定性 FMEA；弹性固定使用 HeuristicSequential、1 h 步长和共同显式故障集。
完整赋值见 \srcpath{src/analysis/counterfactual\_planning.cpp:evaluate\_system}。

对每个存在且有限的维度，基线 $b_d$ 与反事实 $y_d$ 的收益和相对收益为
\begin{align}
 B_d&=b_d-y_d,\\
 s_d&=\max(|b_d|,|y_d|,10^{-9}),\\
 B_d^{rel}&=B_d/s_d.
\end{align}
只有 $B_d>10^{-9}s_d$ 才计为改善，$B_d<-10^{-9}s_d$ 才计为恶化。
同时有改善和恶化为 \texttt{conflict}；否则改善维数至少 2 为 \texttt{compatible}，等于 1
为 \texttt{single\_benefit}，仅有恶化为 \texttt{adverse}，其余为 \texttt{neutral}。
依据为 \srcpath{src/analysis/counterfactual\_planning.cpp:benefit\_between}。

两措施协同只在组合收益和两个单措施收益三者均存在的维度上计算：
\begin{align}
 I_d&=B_{a+b,d}-B_{a,d}-B_{b,d},\\
 s_d^I&=\max(|B_{a+b,d}|,|B_{a,d}|+|B_{b,d}|,10^{-9}),\\
 I_d^{rel}&=I_d/s_d^I.
\end{align}
正负分类仍使用 $10^{-9}s_d^I$；同时有正负为 \texttt{mixed\_interaction}，仅正为
\texttt{synergy}，仅负为 \texttt{redundancy\_or\_conflict}，否则为 \texttt{additive}。
每个单措施和措施对都应用于同一基线的独立副本；措施对按先 $i$ 后 $j$ 的顺序应用。
依据为 \srcpath{src/analysis/counterfactual\_planning.cpp:synergy\_between} 和
\srcpath{src/analysis/counterfactual\_planning.cpp:run\_counterfactual\_planning\_assessment}。

\subsection{覆盖边界}

候选措施生成和五类措施的结构体字段赋值属于工程启发式，不是连续优化数学模型；其完整常数和
分支位于 \srcpath{src/analysis/counterfactual\_planning.cpp:generate\_counterfactual\_measures}。本章未用抽象投资优化模型替代这些赋值。


\section{算法流程}
\subsection{承载力评估主流程}
\begin{figure}[H]\centering
\begin{tikzpicture}[font=\footnotesize,
  box/.style={draw,rounded corners,minimum height=0.85cm,align=center,inner sep=3pt},arr/.style={->,thick}]
  \node[box](agg) at (0,0){母线聚合\\$P,P_G$,DR,$P_{\text{ESS}}$};
  \node[box](bfs) at (3.4,0){供电区 BFS\\同压邻接+变压器边界};
  \node[box](sd) at (6.8,0){设备级 $S_d$\\$\beta$ 自动/手动};
  \node[box](cap) at (10.2,0){$C_{d1},C_{d2}$\\变压器分级};
  \node[box](area) at (3.4,-1.9){区域聚合\\县域 $C_{s1}/C_{s2}$};
  \node[box](sub) at (6.8,-1.9){下级服从上级\\330→10\,kV};
  \node[box](ver) at (10.2,-1.9){工程校验(可选)\\潮流/短路/谐波};
  \draw[arr](agg)--(bfs);\draw[arr](bfs)--(sd);\draw[arr](sd)--(cap);
  \draw[arr](cap.south)|-(area.east);\draw[arr](area)--(sub);\draw[arr](sub)--(ver);
\end{tikzpicture}
\caption{承载力评估流程（工程校验默认关闭）}
\end{figure}

\subsection{薄弱环节筛查伪码}
\begin{quote}\footnotesize\ttfamily
for each entity i:\\
\hspace*{1em}collect available dims K\_i; skip if |K\_i| < min\_evidence\\
\hspace*{1em}p\_ik = max(x\_ik,0)/tau\_k for k in K\_i\\
\hspace*{1em}severity = max p\_ik; consensus = min p\_ik; disagreement = severity-consensus\\
\hspace*{1em}classify(compatible|conflict|single|watch|insufficient)\\
compute Pareto weak frontier within equal (component-group, coverage)\\
sort lexicographic(Pareto, -severity, -consensus, key); take top\_k
\end{quote}

\subsection{反事实评估}
\begin{quote}\footnotesize\ttfamily
base = metrics(system)\\
for each measure j (enabled families, up to max\_measures):\\
\hspace*{1em}sys\_j = apply(copy(system), j); m\_j = metrics(sys\_j); dm\_j = base - m\_j\\
for each pair (j,l): interaction = metrics(apply(apply(copy,j),l)) vs dm\_j+dm\_l\\
report baseline, measures[], interactions[]
\end{quote}

\section{工程假设与数值稳定性分析}
\subsection{工程假设}
承载力假设供电区可由同压支路邻接明确划分、变压器为区域边界、并联变压器容量线性叠加；
薄弱环节假设四维证据已归一化到可比阈值、缺失即真缺失（非零）；反事实假设措施在副本上
的物理施加不破坏引擎可解性。三者均为\textbf{稳态/统计后处理}，不含时步积分。

\subsection{数值稳定性}
承载力为代数闭式求解，唯一数值风险是 $\tau_{\max}\to0$ 时的除零，实现按下界钳位规避；
分级用区间下限，避免边界抖动导致等级翻转。薄弱环节的 $\max/\min$ 聚合对量纲敏感，
故先做阈值归一化 $p_{ik}=x_{ik}/\tau_k$ 使各维可比；非补偿聚合天然对单点异常鲁棒
（一个坏维度不会被平均掉，也不会污染 $c_i$）。反事实的稳定性取决于底层引擎——单条措施
不可解时记为无效并跳过，不影响其余措施。

\section{复杂度分析}
设供电区变压器数 $N_T$、母线数 $N$、实体数 $M$、维度数 $D=4$、措施数 $J$。承载力主体为
每台变压器一次下游 BFS，总计 $O(N_T\cdot(N+E))$（$E$ 为支路数），工程校验开启时另加
一次潮流/短路的引擎成本。薄弱环节筛查为 $O(M D)$ 归一化 + $O(M^2 D)$ Pareto 两两比较
（同组内），实践中 $M$ 为受关注实体数、规模可控。反事实为 $O(J)$ 次单措施引擎求解 +
$O(J^2)$ 次两两协同求解，是三者中最重的，故 \fld{max\_measures} 默认限为 5。

\section{异常工况处理}
\begin{itemize}
  \item \textbf{无下游供电区/孤立母线}：该变压器 $S_d$ 记为不可评估，写入 \fld{warnings}；
        孤立/失电母线不计入任何下游区聚合。
  \item \textbf{$\tau_{\max}$ 或容量为零}：按下界钳位，避免除零；结果标注为退化。
  \item \textbf{维度证据缺失}：排除出 $K_i$，$|K_i|$ 低于要求时判“证据不足”，不参与排序。
  \item \textbf{反事实措施不可施加/不可解}：跳过该措施并在 \fld{methodology} 记录，
        不阻断其余措施与协同分析。
  \item \textbf{工程校验引擎失败}：谐波校核为尽力而为（best-effort），失败不阻断潮流/短路
        主校核，仅在 \fld{verification} 标注缺项。
\end{itemize}

\section{与其他模块的接口关系}
承载力工程校验复用潮流（\srcpath{hacdcpf::solve\_power\_flow}）、详细短路批处理
（\srcpath{analysis::run\_short\_circuit\_detailed\_batch}）与谐波潮流
（\srcpath{harmonics::solve\_harmonic\_power\_flow}）；薄弱环节消费经济（市场/时序）、
碳（\srcpath{carbon\_analysis}）、可靠性（\srcpath{reliability}）、弹性
（\srcpath{resilience}）四引擎的输出证据；反事实经 \srcpath{apply\_counterfactual\_measure}
在系统副本上施加措施后调用同一批引擎。HTTP 侧新增
\texttt{POST /api/session/run\_hosting\_capacity}（\srcpath{tests/run\_gui\_server.cpp}），
序列化经 \srcpath{src/io/json\_io.cpp} 的 \fld{cap\_*} 编解码器往返。

\section{测试与验证建议}
注册测试：\srcpath{test\_hosting\_capacity}、\srcpath{test\_multidimensional\_weak\_link}、
\srcpath{test\_counterfactual\_planning}，以及浏览器端 E2E
\srcpath{multidimensional\_weak\_link\_gui\_test}，覆盖承载力公式、四维筛查排序与反事实
措施评估的一致性。建议补充：$\tau_{\max}\to0$、并联变压器 $\beta$ 自动分支、单维/冲突/兼容
分类边界（$h,c,g$ 邻域）、Pareto 覆盖集不一致时的比较拒绝、以及反事实协同项符号
（增益 vs 抵消）的对照断言。

\section{工业级评估指标}
建议以下指标纳入工程验收：承载力分级与人工审查的一致率（绿/黄/红混淆矩阵）、工程校验
开启前后接入建议的翻转率；薄弱环节 Top-$k$ 命中率（与事后实际故障/整改清单对照）、
分类稳定性（阈值 $\pm5\%$ 扰动下等级不变比例）；反事实收益排序与实际投资后验收益的
秩相关（Spearman $\rho$）、协同项预测符号准确率。全部指标应在固定算例族上可复现。

\section{适用范围与局限性}\label{sec:limits}
\begin{itemize}
  \item 承载力：规范 §5/§6/§9 的系统级生产模拟、县域分解与系统—设备协调\textbf{未}实现；
        \fld{enable\_verification} 默认关闭，即默认仅规划筛查；谐波潮流校核为尽力而为；
  \item 薄弱环节：仅做筛查与决策路由，\textbf{不}宣称归因排放/停电量等于反事实投资收益；
        缺失维度不被当作零；Pareto 仅在同组同覆盖集内成立；
  \item 反事实：措施施加于系统副本，经济/可靠性/弹性指标使用选项给定的有界预算与时窗，
        非全时序生产模拟；
  \item 反事实规划暂无独立契约文档，口径以源码与本手册为准。
\end{itemize}

\section{后续改进方向}
\begin{enumerate}
  \item 承载力向规范 §5/§6/§9 扩展：接入 8760\,h 分区数据与县域分解，形成系统—设备协调；
  \item 薄弱环节引入时段间的风险传播（马尔可夫/生存分析）替代当前逐时段独立投影；
  \item 反事实从两两协同扩展到三阶及以上组合，并以代理模型（surrogate）降低 $O(J^2)$ 引擎
        调用成本；
  \item 为反事实规划补齐独立契约文档，统一四维度量的口径与单位声明。
\end{enumerate}

\section{跨模块工业级技术评价基线}

本章规定本手册所述模块进入工程算例、批量研究或生产服务前必须共同满足的评价基线。
具体模型公式、变量、约束和专用算法以正文相应章节为准；本章不扩大源码已经实现的范围。

\subsection{输入、输出、边界条件与数据结构审计}
输入审计应逐项检查有限性、单位、上下界、枚举值和引用完整性。系统对象必须先按所需层级完成
校验；AC 与 DC 母线采用域限定映射，组件稳定 \texttt{index}、向量位置和临时图索引不得混用。
输出审计必须记录单位、索引空间、模型范围、收敛或可行状态、后端、容差、迭代/节点计数和告警。
空系统、空时域、孤岛、重复 ID、零容量、退化约束与不可用可选后端均属于边界测试集。

\subsection{工程假设与数值稳定性分析}
模型使用的平衡、线性化、稳态、独立性、概率分布、时间离散或网络等值假设必须与应用场景匹配。
数值评价至少包含变量缩放、绝对/相对残差、可行性违反、条件数或枢轴质量、步长/阻尼、迭代停滞
和随机种子复现性。若采用正则化、平滑、回退、松弛或近似，应同时报告触发条件、参数值、对主指标
的敏感性以及是否改变模型口径；不得用“已返回结果”替代收敛或有效性证明。

\subsection{复杂度与资源分析}
令 $n_x$、$n_c$、$n_t$、$n_s$、$|V|$、$|E|$ 分别表示变量、约束、时段、场景、节点和边数。
报告应分别给出建模、求解、后处理的时间与内存。图遍历通常为 $O(|V|+|E|)$；逐时段/逐场景
流程至少线性增长为 $O(n_t n_s)$ 次子问题调用；稠密线性代数最坏可达 $O(n_x^3)$，稀疏方法则应
报告结构非零数、填充率和因子分解峰值内存；MILP/NLP 不给出虚假的多项式最坏界，而报告节点数、
间隙、时限和实测扩展曲线。复杂度结论必须基于固定硬件、固定后端和可复现算例族。

\subsection{异常工况处理}
非法输入应在进入求解器前返回结构化错误；不可行、无界、不收敛、时限、内存不足和后端缺失必须
相互区分。部分结果只有在对应 \fld{ValidityFlags}、\fld{model\_scope}、
\fld{model\_limitations} 或等价字段允许时才可使用。异常路径不得静默丢弃 AC/DC 资产、制造平衡
节点、把 NaN 改写为零或把近似解标为全局最优；系统原始对象应保持可恢复和可重试。

\subsection{与其他模块的接口关系}
接口链路遵循“工程语义模型 $\to$ 校验 $\to$ 投影/装配 $\to$ 求解或分析 $\to$ 结果回投影”。
跨模块传递时保留稳定 ID、单位、符号、时间基准和场景身份；消费潮流、OPF、动态或可靠性结果前，
先检查其收敛、覆盖范围和有效性标志。HTTP/JSON/Python 层只做语义保持适配，不重新解释求解结果。

\subsection{测试与验证建议}
最低验证矩阵包括：手算小例、标准算例、边界/异常输入、随机性质测试、算法或后端交叉校核、
序列化往返、稳定 ID 回投影、AC/DC 同号母线、重复运行确定性和规模扩展测试。数值算法还应加入
残差独立重算、雅可比有限差分或自动微分校核；近似模型应与高保真模型比较并给出误差分布，而非
只报告单个平均值。回归报告必须列明构建配置、算例版本、容差、通过/失败/跳过数及未验证范围。

\subsection{工业级评估指标}
\begin{itemize}
  \item \textbf{正确性}：守恒残差、约束最大违反、解析/基准误差、结果回投影一致性；
  \item \textbf{鲁棒性}：收敛率、不可行识别率、异常分类准确率、扰动敏感性与随机复现率；
  \item \textbf{性能}：端到端时延、吞吐量、峰值内存、并行效率与规模增长斜率；
  \item \textbf{可追溯性}：输入模式版本、稳定 ID 覆盖率、模型范围/限制字段完整率、告警可定位率；
  \item \textbf{工程有效性}：与标准、外部引擎或现场数据的偏差，以及误判/漏判对业务指标的影响。
\end{itemize}
指标应预先给出验收阈值和失败处置；未达到阈值时只能标记为实验性或受限适用。

\subsection{适用范围、局限性与后续改进}
本手册只评价当前源码覆盖的模型、后端和数据结构，不对未建模物理过程、未安装求解器或未执行的
外部验证作保证。后续改进应按“缺口 $\to$ 可量化指标 $\to$ 源码/测试变更 $\to$ 独立验证”闭环，
优先消除会造成错误结论的语义缺口，其次提升数值鲁棒性与性能，最后扩展模型范围。


\end{document}
